\chapter{Representación de las rutas y función de costo}
\label{cap:repr}

\section{Cómo representamos una solución}
Una de las ideas que más se repite en la investigación previa es que la forma de representar el
problema influye mucho en qué tan bien funciona el algoritmo. En nuestro caso, toda solución
candidata es una \textbf{permutación} de los vértices, es decir, una lista que contiene cada
vértice una sola vez. Una permutación puede ser o no un ciclo hamiltoniano, según las aristas que
falten, pero nunca es una ruta ``ilegal'' con vértices repetidos u olvidados.

Cada algoritmo se asegura de mantener esta regla a su manera:
\begin{itemize}
  \item \textbf{GA, SA y ABC} trabajan directamente con permutaciones, y sus operaciones
        (cruzar, intercambiar dos vértices, invertir un tramo) siempre devuelven otra permutación.
  \item \textbf{ACO} construye las rutas avanzando solo por aristas que existen. Si una hormiga
        se queda atorada, se completa la ruta con los vértices que le faltan.
  \item \textbf{DE y PSO} están pensados para trabajar con números reales, no con listas. Por eso
        usamos \emph{random keys}: cada solución es un vector de $n$ números entre 0 y 1, uno por
        vértice, y la ruta se obtiene ordenando los vértices según su número. Por ejemplo,
        $(0.7,\ 0.1,\ 0.4)$ da la ruta $(2,3,1)$. Como ordenar siempre da una permutación, estos
        algoritmos tampoco pueden generar rutas inválidas (\texttt{src/representacion.py}).
\end{itemize}

Al final de los experimentos verificamos que las 2\,112 rutas guardadas fueran permutaciones
válidas, y todas lo fueron (sección~\ref{sec:verificacion}).

\section{Función de costo}
Para que los algoritmos sepan si van bien, necesitan una forma de calificar cada ruta. La que
usamos es muy simple: \textbf{contar cuántas conexiones de la ruta no existen en el grafo},
incluyendo la conexión entre el último vértice y el primero.
\[
f(\pi)=\text{número de pares consecutivos }(v_i,v_{i+1})\text{ de la ruta que no son arista del grafo}.
\]
Si $f(\pi)=0$, la ruta es un ciclo hamiltoniano, así que el objetivo de todos los algoritmos es
llevar este número a cero. En el ejemplo de la sección anterior, la ruta $(1,3,2,4,5)$ tiene
$f=1$ porque solo le falta la arista $2\!-\!4$.

Escogimos esta función por tres razones:
\begin{enumerate}
  \item Es fácil de entender: el valor dice directamente cuántas aristas faltan.
  \item Es la misma para todos los algoritmos, así que la comparación es justa.
  \item Se puede actualizar muy rápido. Cuando un algoritmo intercambia dos vértices o invierte
        un tramo de la ruta, solo cambian unas cuantas conexiones, así que no hace falta revisar
        la ruta completa: basta con revisar las conexiones que cambiaron. Comprobamos que este
        cálculo rápido coincide con el cálculo completo en 25\,440 casos de prueba
        (\texttt{python -m src.fitness}).
\end{enumerate}

Su desventaja es que muchas rutas distintas tienen el mismo costo. Esto crea zonas ``planas''
donde el algoritmo no sabe hacia dónde moverse. Es una de las razones por las que algunos
algoritmos se estancan, como veremos en los resultados.

\section{Un movimiento guiado por el grafo}
Además de los movimientos al azar, programamos un movimiento ``inteligente'' que llamamos
\textbf{inversión guiada} (\texttt{src/operadores.py}). Funciona así:
\begin{enumerate}
  \item Busca en la ruta una conexión que no existe, por ejemplo entre $u$ y el vértice que le
        sigue.
  \item Elige un vecino real $w$ de $u$ en el grafo.
  \item Invierte el tramo de la ruta que está entre ellos, de modo que $w$ quede justo después
        de $u$.
\end{enumerate}
Así, en vez de mover vértices al azar, el algoritmo intenta reparar una conexión rota usando una
arista que sí existe. Lo usa principalmente ABC, y como veremos, hace una gran diferencia.
